% titel: Blatt 6, A9: Polare Bindungen und Dipolmoment

\subsection*{Lösung}
\begin{center}
\begin{tabular}{@{}l l l@{}}
  Molekül & polare Bindungen? & permanentes Dipolmoment? \\
  \hline
  $\mathrm{I_2}$      & nein (homoatomar)  & nein \\
  $\mathrm{HCl}$      & ja, stark polar     & ja \\
  $\mathrm{CH_3OH}$   & ja ($\mathrm{C-O}$, $\mathrm{O-H}$) & ja \\
  $\mathrm{C_6H_6}$   & praktisch unpolare C--H-Bindung & nein (symmetrisch)
  \\
  $\mathrm{H_2O}$     & ja                  & ja \\
  $\mathrm{CO_2}$     & ja ($\mathrm{C=O}$)& nein (linear) \\
  $\mathrm{CH_2Cl_2}$ & ja                  & ja (asymmetrisch) \\
  $\mathrm{CHCl_3}$   & ja                  & ja (asymmetrisch) \\
  $\mathrm{CCl_4}$    & ja ($\mathrm{C-Cl}$)& nein (symmetrisch-tetraedrisch) \\
  $\mathrm{CS_2}$     & ja ($\mathrm{C=S}$)& nein (linear) \\
  $\mathrm{SF_4}$     & ja                  & ja (see-saw, asymmetrisch) \\
\end{tabular}
\end{center}

\subsection*{Begründung}
$\mathrm{CO_2}$, $\mathrm{CS_2}$ (linear), $\mathrm{CCl_4}$ und $\mathrm{C_6H_6}$
haben zwar polare Einzelbindungen, aber die Vektorsumme verschwindet aus
Symmetrie. $\mathrm{CH_2Cl_2}$ und $\mathrm{CHCl_3}$ besitzen asymmetrische
Ligandenverteilung $\to$ Nettodipol. $\mathrm{SF_4}$ (see-saw) kompensiert
nicht vollständig. $\mathrm{CH_3OH}$, $\mathrm{HCl}$, $\mathrm{H_2O}$
sind unsymmetrisch und damit polar.

\begin{ergebnis}
Dipolmoment: ja bei $\mathrm{HCl}$, $\mathrm{CH_3OH}$, $\mathrm{H_2O}$,
$\mathrm{CH_2Cl_2}$, $\mathrm{CHCl_3}$, $\mathrm{SF_4}$; nein bei
$\mathrm{I_2}$, $\mathrm{C_6H_6}$, $\mathrm{CO_2}$, $\mathrm{CCl_4}$,
$\mathrm{CS_2}$ (Symmetrie). Alle bis auf $\mathrm{I_2}$ und
$\mathrm{C_6H_6}$ enthalten polare Bindungen.
\end{ergebnis}

\begin{kontrolle}
Gruppierung: symmetrische Geometrien (linear/tetraedrisch, planare
Ringe) $\to$ kein Dipol; Asymmetrie $\to$ Dipol. Bei $\mathrm{C_6H_6}$ ist die
C--H-EN-Differenz klein und die Geometrie hochsymmetrisch --- übereinstimmend
mit dem bekannten experimentellen Befund (Benzol apolar).
\end{kontrolle}

\begin{merksatz}
Polarität einer Bindung: EN-Differenz. Dipolmoment: Symmetrie entscheidet ---
nicht jede polare Bindung ergibt ein polares Molekül. Lineare und
tetraedrische symmetrische Moleküle kompensieren ihre Dipole.
\end{merksatz}
